\documentclass[11pt]{article}
\usepackage[utf8]{inputenc}
\usepackage[T1]{fontenc}
\usepackage{amsmath, amssymb, amsthm}
\usepackage{geometry}
\usepackage{booktabs}
\usepackage{hyperref}

\geometry{margin=1in}

\title{Foundations of Kähler and Hyperkähler Geometry}
\author{Technical Overview}
\date{\today}

\begin{document}

\maketitle

\section{Introduction}
In the study of differential geometry, Kähler and Hyperkähler manifolds represent a critical intersection of Riemannian, complex, and symplectic structures. These manifolds are central to modern theoretical physics, particularly in string theory and supersymmetric field theories.

\section{Kähler Manifolds}
A Kähler manifold is a manifold $M$ equipped with a Riemannian metric $g$ and a complex structure $J$ that are compatible. The fundamental 2-form (or Kähler form) is defined by:
\begin{equation}
    \omega(X, Y) = g(JX, Y)
\end{equation}

The manifold is Kähler if and only if $\omega$ is closed ($d\omega = 0$). This condition implies that the holonomy group of the metric is contained within the unitary group $U(n)$ \cite{huybrechts}.

\subsection{Key Properties}
\begin{itemize}
    \item \textbf{Kähler Potential:} Locally, the metric can be derived from a scalar potential $\Phi$ such that $g_{i\bar{j}} = \partial_i \partial_{\bar{j}} \Phi$.
    \item \textbf{Hodge Decomposition:} The cohomology of a compact Kähler manifold admits a decomposition $H^k(M, \mathbb{C}) = \bigoplus_{p+q=k} H^{p,q}(M)$.
\end{itemize}

\section{Hyperkähler Manifolds}
A Hyperkähler manifold is a Riemannian manifold $(M, g)$ that admits three distinct complex structures $I, J, K$ satisfying the quaternionic relations:
\begin{equation}
    I^2 = J^2 = K^2 = IJK = -1
\end{equation}

For such a manifold, the metric is Kähler with respect to all three structures simultaneously. This forces the holonomy group to be a subgroup of the symplectic group $Sp(n)$, which implies that the manifold is Ricci-flat ($R_{ij} = 0$) \cite{beauville}.

\section{Comparison of Structures}
The following table summarizes the structural differences between these two classes of manifolds.

\begin{table}[h]
\centering
\begin{tabular}{@{}lll@{}}
\toprule
\textbf{Property} & \textbf{Kähler} & \textbf{Hyperkähler} \\ \midrule
Complex Structures & 1 ($J$) & 3 ($I, J, K$) \\
Holonomy Group & $U(n)$ & $Sp(n)$ \\
Curvature Constraint & None (General) & Ricci-Flat \\
Real Dimension & $2n$ & $4n$ \\ \bottomrule
\end{tabular}
\caption{Comparative analysis of Kähler and Hyperkähler properties.}
\end{table}

\begin{thebibliography}{9}

\bibitem{huybrechts}
D. Huybrechts, \textit{Complex Geometry: An Introduction}, Universitext, Springer-Verlag, Berlin, 2005.

\bibitem{beauville}
A. Beauville, \textit{Variétés Kählériennes dont la première classe de Chern est nulle}, J. Differential Geom. 18 (1983), no. 4, 755--782.

\bibitem{moroianu}
A. Moroianu, \textit{Lectures on Kähler Geometry}, London Mathematical Society Student Texts, Cambridge University Press, 2007.

\bibitem{hitchin}
N. J. Hitchin, A. Karlhede, U. Lindström and M. Roček, \textit{Hyperkähler metrics and supersymmetry}, Comm. Math. Phys. 108 (1987), no. 4, 535--589.

\end{thebibliography}

\end{document}
